\documentclass[12pt]{article}

\usepackage{bbm}
\usepackage{fancyhdr}
\usepackage{titling}
\usepackage{enumitem}
\usepackage{amsmath}
\usepackage{dsfont}
\usepackage{amsfonts}
\usepackage{amssymb}
\usepackage{amsthm}
\usepackage{geometry}
\usepackage{tikz}
\usepackage{algorithm}
\usepackage{algpseudocode}
\geometry{verbose,tmargin=1in,bmargin=1in,lmargin=1in,rmargin=1in}
\renewcommand{\vec}[1]{\mathbf{#1}}
\newcommand{\dec}{\text{dec}}
\newcommand{\divi}{\mathbin{\textnormal{div}}}
\title{Homework 1}
\author{Satvik Prasad}
\date{}
\newtheorem{lemma}{Lemma}[section]
\fancypagestyle{plain}{%
  \fancyhf{}
  \fancyhead[L]{\theauthor}
  \fancyhead[R]{\thedate}
}
\pagestyle{plain}
\makeatletter
\def\@maketitle{%
  \newpage
  \null
  \vskip 1em%
  \begin{center}%
    \let \footnote \thanks
    {\LARGE \@title \par}%
    \vskip 1em%
  \end{center}%
  \par
\vskip 1em}
\makeatother
\usepackage{cmbright}

\newcommand{\mat}[1]{\mathbf{#1}}

% ---------- Probability theory ----------
% Core operators (optional argument-free; use as \E{X}, \Var{X}, \Cov{X}{Y})
\DeclareMathOperator{\EOp}{\mathbb{E}}
\DeclareMathOperator{\VarOp}{Var}
\DeclareMathOperator{\CovOp}{Cov}
\DeclareMathOperator{\CorrOp}{Corr}

\renewcommand{\P}[1]{\mathbb{P}\left(#1\right)}        % \P{A}
\newcommand{\E}[1]{\EOp\left[#1\right]}                % \E{X}
\newcommand{\Var}[1]{\VarOp\left(#1\right)}            % \Var{X}
\newcommand{\Cov}[2]{\CovOp\left(#1,\,#2\right)}       % \Cov{X}{Y}
\newcommand{\Corr}[2]{\CorrOp\left(#1,\,#2\right)}     % \Corr{X}{Y}
\newcommand{\SD}[1]{\operatorname{SD}\left(#1\right)}  % \SD{X}

% Subscripted versions, e.g. \Esub{\theta}{X}, \Psub{\theta}{A}
\newcommand{\Esub}[2]{\EOp_{#1}\left[#2\right]}
\newcommand{\Psub}[2]{\mathbb{P}_{#1}\left(#2\right)}

% Conditioning, indicators, independence
\newcommand{\given}{\,\middle|\,}                      % \E{X \given Y}
\newcommand{\ind}[1]{\mathbbm{1}\left\{#1\right\}}     % \ind{X > 0}
\newcommand{\indep}{\perp\!\!\!\perp}                  % X \indep Y
\newcommand{\iid}{\overset{\text{iid}}{\sim}}          % X_i \iid \Normal(0,1)

% Convergence
\newcommand{\convp}{\xrightarrow{\;p\;}}               % in probability
\newcommand{\convd}{\xrightarrow{\;d\;}}               % in distribution
\newcommand{\convas}{\xrightarrow{\text{a.s.}}}        % almost surely
\newcommand{\convLp}[1]{\xrightarrow{L^{#1}}}          % in L^p

% Probability spaces
\newcommand{\F}{\mathcal{F}}                           % sigma-algebra
\newcommand{\B}{\mathcal{B}}                           % Borel sigma-algebra
\newcommand{\pspace}{(\Omega, \F, \mathbb{P})}

% Common distributions
\newcommand{\Normal}{\mathcal{N}}
\newcommand{\Bern}{\operatorname{Bernoulli}}
\newcommand{\Bin}{\operatorname{Binomial}}
\newcommand{\Pois}{\operatorname{Poisson}}
\newcommand{\Unif}{\operatorname{Uniform}}
\newcommand{\Exp}{\operatorname{Exp}}
\newcommand{\Geom}{\operatorname{Geometric}}
\newcommand{\Gam}{\operatorname{Gamma}}
\newcommand{\Beta}{\operatorname{Beta}}

% Number sets
\newcommand{\R}{\mathbb{R}}
\newcommand{\N}{\mathbb{N}}
\newcommand{\Z}{\mathbb{Z}}


\begin{document}
\begin{enumerate}
  \item \textbf{Green Circle, 1} Construct the map $f: \{a_1, \ldots,
    a_{n + 1}\} \to \{1, \ldots, n\}$ such that
    $$
    f(2k - 1) = k, \qquad f(2k) = k, \qquad \forall k \in \{1, \ldots, n\}.
    $$
    By the pigeonhole principle, this map cannot be injective, and
    there must be some $i \neq j$ such that $f(a_i) = f(a_j).$ Since
    $a_i, a_j$ distinct, (assuming WLOG that $a_i < a_j$), this means
    $$
    a_j = a_i + 1.
    $$
    Now, for the sake of contradiction assume that $a_j, a_i$ aren't
    relatively prime. Then, there exists a prime factor $p,$ and
    positive integers $M, N$ with $N > M$ such that
    $$
    a_i = pM, \ a_j = pN \Rightarrow pN = pM + 1,
    $$
    meaning $p(N - M) = 1.$ Since $p > 1,$ this means $M< N < M + 1,$
    which is impossible for an integer $N.$ Thus, this is a
    contradiction, meaning $a_i, a_j$ are relatively prime,
    concluding our proof.

    We show that $n + 1$ cannot be replaced by $n$ by considering the
    case where $n = 2.$ Then, we can choose $a_1 = 2, a_2 = 4$ such
    that $1 \leq a_1, a_2 \leq 2(2),$ and $a_1, a_2$ are not
    relatively prime. Hence, $n + 1$ is the smallest such integer
    such that this holds in the general case.

  \item \textbf{Green Circle, 2} Notice that
    $$
    \dfrac{a_i - a_j}{(1 +a_i)(1+a_j)} = \dfrac{1}{1+a_j} - \dfrac{1}{1+a_i}.
    $$
    Since $a_k \in \mathbb R^{+},$ we know
    $$
    a_j > 0 \Rightarrow 1 + a_j > 1 \Rightarrow \dfrac{1}{1+a_j} < 1.
    $$
    Additionally, since $a_j > 0,$ we know
    $$
    \dfrac{1}{1+a_j} > 0.
    $$

    Then, construct the map $f: \{a_1, \ldots, a_{n +1}\} \to \{1,
    \ldots, n\}$ defined by
    $$
    f(x) = \left\lceil \dfrac{n}{1+x} \right\rceil.
    $$

    By the pigeonhole principle, there must be $i \neq j$ such that
    $\lceil n/(1 + a_i) \rceil = \lceil n/(1 + a_j) \rceil.$ In otherwords,
    $$
    k - 1< n/(1 + a_j) \leq k, \qquad k - 1< n/(1 + a_j) \leq k
    $$
    for some integer $1 \leq k \leq n.$

    Manipulating, we get
    $$
    \frac{k - 1}{n}< 1/(1 + a_j) \leq \frac{k}{n}, \qquad \frac{k -
    1}{n}< 1/(1 + a_j) \leq \frac{k}{n},
    $$
    meaning
    $$
    -\frac{1}{n}< \dfrac{1}{1+a_i} - \dfrac{1}{1 +a_j} < \frac{1}{n}
    \Rightarrow \left|\dfrac{a_i - a_j}{(1+a_i)(1+a_j)}\right| < \dfrac{1}{n}.
    $$
    Thus, (assuming WLOG $a_i > a_j$), we have
    $$
    0 < \dfrac{a_i - a_j}{(1+a_i)(1+a_j)} \leq \dfrac{1}{n}.
    $$

  \item \textbf{Green Circle, 4} Represent this as a graph with
    vertices $v_k$ for each person, and edge $[v_i, v_j]$ if $i$ is
    an acquaintance of $j.$ Then, we know
    $$
    \sum_{k = 1}^{n} \deg(v) = 2|E|,
    $$
    meaning the sum of the degrees of vertices must be even.

    Proceed via induction. For our base case, when $n = 2,$ two
    vertices can either both have degree $0,$ or both have degree
    $1,$ meaning there must be two vertices with the same degree.

    Assume for $n - 1 \geq 1$ vertices in a bidirectional graph, two
    must have the same degree. Then, consider a graph with degree $n.$

    A vertex can have degree $0, 1, \ldots, n- 1.$ We split this into two cases:
    \begin{itemize}
      \item If there is no vertex with degree $0,$ then there are $n
        - 1$ possible degrees, and $n$ vertices, meaning two vertices
        must have the same degree.
      \item If there is a vertex of degree $0,$ $v_i,$ the remaining
        vertices form a graph with $n - 1$ vertices. By our inductive
        hypothesis, two vertices must have the same degree in this graph.
    \end{itemize}

    Therefore, by induction, a bidirectional graph must have two
    vertices with the same degree. By our bijection, there must be
    two people with the same number of acquaintances.

  \item \textbf{Blue Square, 6} Some non-empty consecutive block is a
    perfect square if and only if every prime $p$ that divides $a_{r
    + 1}, \ldots, a_{s}$ appears an even number of times.

    Write $a_i = p_1^{r_{1i}} p_2^{r_2i} \ldots p_k^{r_ki}$ Then, for
    a consecutive block to be a perfect square, we need consecutive
    indices $r + 1, \ldots, s$ such that
    $$
    \sum_{i = r+1}^{s} r_{ji} \equiv 0 \mod 2
    $$
    for all integers $1 \leq j \leq k.$ Let $\eta : \{0, \ldots,
    2^k\} \to \{0, 1\}^k$ be defined by
    $$
    \eta_j(s) = \left(\sum_{i = 1}^{s} r_{ji}\right) \bmod 2, \qquad
    \eta(s) = (\eta_1, \ldots, \eta_k).
    $$
    Notice that $|\{a : a = \eta(s)\}| = 2^k,$ since each entry can
    either be $0$ or $1.$ Thus, by the pigeonhole principle there
    must be two unique $s_1, s_2$ (representing prefixes of $a_1,
    \ldots, a_{2^k}$) such that $\eta(s_1) = \eta(s_2).$ WLOG, assume
    $s_1 < s_2.$ Then, consider
    $$
    a_{s_1 +1}, \ldots a_{s_2}.
    $$
    Notice that
    \begin{align*}
      \sum_{i = s_1 + 1}^{s_2} r_{ji} = \eta_j(s_2) - \eta_j(s_1)
      \equiv 0 \mod 2,
    \end{align*}
    meaning this consecutive, non-empty block is a perfect square.

    Thus, by the pigeonhole principle there must be at least one
    consecutive non-empty block that is a perfect square.

  \item \textbf{Blue Square, 11}
    \begin{enumerate}
      \item We want to show that for any irrational $\alpha,$ there
        exists positive integers $p, q$ such that
        $$
        \left|\alpha - \dfrac{p}{q}\right| < \dfrac{1}{Nq}.
        $$

        In other words, we need
        $$
        |q\alpha - p| < \dfrac{|q|}{Nq}.
        $$
        Then, we have
        $$
        |q \alpha - p| < \dfrac{1}{N},
        $$

        Consider the mapping $f : \{0, \ldots, N\} \to \{1, \ldots,
        N\}$ given by
        $$
        f(k) = \lceil N \{k \alpha\} \rceil.
        $$

        By the pigeonhole principle, there must be two $i \neq j$
        such that $f(i) = f(j).$ WLOG assume $i < j$ Then, we have
        $$
        \lceil N\{j \alpha\} \rceil = \lceil N\{i \alpha \}\rceil,
        $$
        meaning
        $$
        |N\{j \alpha\}-N\{i \alpha\}| < 1,
        $$
        an exclusive bound since $\alpha$ irrational. This allows us to conclude
        $$
        | (j-i) \alpha - (\lfloor j \alpha \rfloor - \lfloor i \alpha
        \rfloor)| < \dfrac{1}{N}.
        $$
        Letting $1 \leq q = j - i \leq N,$ and $q = \lfloor j \alpha
        \rfloor - \lfloor i \alpha \rfloor,$ we get
        $$
        |q \alpha - p| < \dfrac{1}{N} \Rightarrow \left|\alpha -
        \dfrac{p}{q}\right| < \dfrac{1}{qN}.
        $$

      \item Assume the set of $p, q$ satisfying
        $$
        \left|\alpha  - \dfrac{p}{q}\right| < \dfrac{1}{Nq}
        $$
        for any $N$ is finite. Then, consider the $(p', q')$ tuple
        that minimizes
        $$
        (p', q') = \text{argmin}_{p, q} \left|q\alpha - p\right|.
        $$

        Choose the smallest integer $N"$ such that $$N' >
        \dfrac{1}{|q'\alpha - p'|}.$$ This satisfies
        $$
        \dfrac{1}{N} < \left|q'\alpha - p'\right|,
        $$
        meaning for all $p, q,$
        $$
        |q\alpha - p| > \dfrac{1}{N'} \Rightarrow \left|\alpha -
        \dfrac{p}{q}\right| > \dfrac{1}{Nq},
        $$
        which contradicts the existence of valid $p, q$ for any $N.$

        Therefore, the set of $p, q$ must be infinite, meaning there
        are infinitely many $p, q$ such that
        $$
        \left | \alpha - \dfrac{p}{q} \right| < \dfrac{1}{Nq}.
        $$
        Since $N \geq q,$ $1/ Nq \leq 1/q^2,$ and thus there are
        infinitely many $p, q$ such that
        $$
        \left|\alpha - \dfrac{p}{q}\right| < \dfrac{1}{q^2}.
        $$
    \end{enumerate}

  \item \textbf{Black Diamond, 12} Insight: we need to make the same
    argument for a single $q$ over multiple $p,$ implying we need the
    same $j, i$ over all $\alpha_1, \ldots, \alpha_d.$ Hence, we need
    to show that there exists some $j, i$ such that the largest
    pairwise distance between corresponding elements in the two sets
    $$
    (\{j \alpha_1\}, \{j\alpha_2\}, \ldots, \{j \alpha_d\}), (\{i
    \alpha_1\}, \{i \alpha_2\} \ldots, \{i \alpha_d\})
    $$
    is $1/N.$

    This corresponds exactly to the two points (treating LHS and RHS
    as a vector) being in the same unit cube.

    Thus, consider the $N^d + 1$ points
    $$
    \vec p_k = (\{k\alpha_1\}, \{k\alpha_2\}, \ldots, \{k\alpha_d\})
    \in \mathbb [0, 1]^{d}, \qquad 0 \leq k \leq N^d.
    $$
    Split $[0, 1]^d$ into $N^d$ cubes of side length $1 /N.$ Then, by
    the pigeonhole principle there must be two indices $i, j$ (WLOG
    assume $j > i$) such that
    $$
    \vec p_i, \vec p_j
    $$
    are in the same unit cube.

    Thus,
    $$
    |\{j \alpha_k\} - \{i\alpha_k\}| < \dfrac{1}{N}
    $$
    for all $1 \leq k \leq d,$ (not an equality because $\alpha_k$ is
    irrational). Equivalently,
    $$
    |(j - i) \alpha_k - (\lfloor j\alpha_k\rfloor - \lfloor i
    \alpha_k\rfloor)| < 1/N.
    $$

    Let $j - i = q,$ and $\lfloor j\alpha_k\rfloor - \lfloor
    i\alpha_k\rfloor =p_k$ for some $k.$ Then, we have
    $$
    |q\alpha_k - p_k| < \dfrac{1}{N}
    $$
    for $1 \leq q \leq N^d$ and all $1 \leq k \leq d$ as required.

    Dividing both sides by $q,$ we can conclude
    $$
    \left|\alpha_k - \dfrac{p_i}{q_i}\right| < \dfrac{1}{qN}, \qquad
    1 \leq i \leq d.
    $$
\end{enumerate}

\subsection{Reflection}
For Blue Square 6, initially looked at all possible contiguous
sequences, but realised here a collision didn't really help if we
define buckets on prime multiplicity parity, because you're not
guaranteed to have something in the bucket corresponding to even
multiplicities. Instead, saw that if we have a collision in the same
bucket, the division of the larger number with the smaller number
will have prime multiplicities that are all even. From there, a
contiguous sequence can be represented as a larger prefix divided by
a smaller prefix, so the solution fell out.

\end{document}
